\scanpage{27}
\mspage{-4}
\diagram{\begin{tikzcd}[column sep=large,row sep=large]
G & & G^p \arrow[ll,"\mu"'] \arrow[dl,"i"'] & & G \arrow[ll,"\delta"'] \arrow[ddll,"F"'] \\
& \Sigma^pG \arrow[ul,"a"] & & & \\
& & G^{\deleted{\unsure{p}}\inserted{(p)}} \arrow[ul,"r"'] & &
\end{tikzcd}}

où $\delta$ est l'application diagonale et $\mu$ l'application $(g_1,\ldots,g_p)\dashrightarrow g_1+\cdots+g_p$; $\Sigma^pG$ est le produit symétrique de G (voir Serre, loc. cit.) ; $(i,a)$ est la décomposition canonique de $\mu$ ; quant à la décomposition $i.\delta=r.F$, elle est définie de façon canonique pour toute variété algébrique (voir Serre, loc.cit.).

Le diagramme ci-dessus est donc défini pour tout groupe algébrique commutatif, affine ou non. On posera donc $V=a.r$ et on aura $V.F=p$.
\mspage{26}
